\begin{center}\hypertarget{native-228-article165}{}\pdfbookmark[2]{165.}{article-165}165.\end{center}
The following appears to be the simplest method of fitting an already approximately known orbit to four observations. Let $x$, $y$ be the approximate distances of the \SourceLeaf{229}celestial body from the Earth at the two complete observations. From these distances compute the corresponding heliocentric positions and hence the elements themselves. Then use these elements to compute the geocentric longitudes or right ascensions for the two remaining observations. If these happen to agree with the observed values, the elements need no further correction. Otherwise, note the differences $X$, $Y$ and repeat the same calculation twice, changing the values of $x$, $y$ slightly. This will give three sets of values of the quantities $x$, $y$ and the differences $X$, $Y$. By following the directions in Article 120, we derive from these the corrected values of $x$, $y$, to which the values $X=0$, $Y=0$ correspond. A similar calculation based on this fourth set will yield elements that duly represent all four observations.

If a choice is possible, it is preferable to retain as complete those observations from which the orientation of the orbit can be determined with the greatest precision, hence the first and last observations whenever they cover a heliocentric motion of $90$ degrees or less. If the observations are of unequal precision, however, set aside the latitudes or declinations of those you suspect to be less accurate.
